\documentclass[12pt]{article}
\usepackage[utf8]{inputenc}
\usepackage[spanish]{babel}
\usepackage {setspace}
\onehalfspacing
\setlength{\parskip}{1em}
%------------------------------------------------


\title{PRIMERA ENTREGA DEL PROYECTO PROTOTIPICO}
\author{Cesar Ivan Rodriguez Navarro \and 
	Victor Manuel De La Cruz \and 
	Frida Yosselin Zaragoza }
\begin{document}

\maketitle 

\begin{abstract}
Plantación De Cebolla en Chihuahua
\end{abstract}

\section{INDICE}
\subsection{Región, Sembradío e Información general}
\subsection{Identificación  De El Problema}
\subsection{Justificación Del Problema}
\subsection{Programación}
\subsection{Ciencia de Datos}
\subsection{Geometría Analítica}
\subsection{Matemáticas Discretas}
\subsection{Administración Empresarial}
\subsection{Perspectiva de Genero}
\subsection{Posibles Soluciones Tempranas}
\subsection{FUENTES CONSULTADAS}
\newpage
\begin{center}
	{\huge \textbf{Región, Sembradío e Información general} }
\end{center}
\vspace{1cm} 
{\large %
Para esta sección escogimos la región de chihuahua, ya que es la región que mas aporta a la producción de cebolla, dedicando prácticamente el 21\% -21.6\% de la producción de todo México, rondando casi las 300,000 toneladas por cosecha.

Es de los estados mas importante, igual que Zacatecas, pero como este es el que genera mas nos concentramos ya que una perdida de 300,00toneladas, suena poco ante las estadísticas, pero si es un gran numero y perdidas tanto en venta como en producción y ganancia de quienes invierten y siembran.

También vimos que el estado de Chihuahua ha sido de los mas afectados , marcado por la CONAGUA como "sequia severa" y buscando la información ha sido bastante impactante el cambio climático que ha pasado durante estos años, alrededor del 2020 al 2025

Justamente la cebolla tarda 4 meses en germinar lo que puede atraer varios problemas , germinan en un 1/3 del año , su termino de cosecha tardan alrededor de los 4 meses , y puede variar su tamaño aunque lleguen a ser semillas de calidad , estas variando entre 4 clases de cebolla variando por su tamaño , abarcando desde los cebollines (estos se encuentran en las taquerias normalmente)\ hasta las cebollas tamaño jumbo las cuales llegan a pesar casi 1kg por pieza, la cual puede ser tanto buena como mala noticia ya que la ganancia variando entre precios y aveces clase de cebolla , sea la blanca que todos conocemos , la cebolla morada o el cebollin de las taquerias.
\newpage
\begin{center}
	{\huge \textbf{Identificación  De El Problema}}
\end{center}
\vspace{1cm} 
{\large %
	Justamente la problemática la vemos desde la parte del cambio climático que se ha generado en estos años recientes, tomando desde el año 2020 hasta el actual, junto con la sobrante lluvia y calores que rondan los 32° pueden dañar severamente el plantío.
	
	Investigando como se siembra y todo el proceso junto con experiencia (uno de nuestros compañeros cosecha cebolla en casa) un repentino cambio de clima tan brusco puede dañar y/o causar enfermedades a las plantaciones, hay varias enfermedades, pero las principales son:
	
	Mancha purpura, Mildiu Lanoso, Raíz Rosada, Podredumbre del Cuello, Fusuariosis y Moho negro 
	Estas enfermedades afectan bastante, tanto como el rendimiento y deterioro de raíces de la misma cosecha, en el sembradío normalmente antes de plantar se “esteriliza” el suelo echando vitaminas, fungicidas y pesticidas antes de la siembra para matar el hongo y varias enfermedades. 
	Esto junto que con la sección de chihuahua ha sido la mas arrasada ante estos cambios climáticos, ha sido una de las regiones con el calentamiento más acelerado junto con sequías severas 
	
	DATOS PARA FUTURAS GRÁFICAS:
	2019- este fue un año de transición, con el aumento de las temperaturas máximas al promedio, en este año seguía algo moderado el estado 
	2020- Aquí hay un cambio relevante, hubo olas de calor intenso, además de un año extremadamente seco, sequías severas y los niveles de las presas bajaron alrededor del 25-30%.
	2021- aquí hubo un cambio totalmente radical (y preocupante a mi opinión), con nevadas históricas en febrero y montón fuerte en verano, aunque aquí hubo un poco recuperación parcial de acuíferos
	2022-verano con lluvias y récord, este año fue caótico por varias inundaciones y sobre abasto de las presas al 100%
	2023- este fue el año mas caluroso registrado a nivel global “ebullición global” aquí hubo un cambio fuerte, regreso de sequías extremas y calor persistente en septiembre 
	2024- Aquí se rompió la gráfica superando las temperaturas máximas de 40 grados en varias regiones, la lluvia se prolongo hasta mediados de año 
	2025-aquí hubo un invierno extremo (como ejemplo hubo un -23 grados en  Majalca ) seguido de un verano muy caluroso , se recuperaron las presas , un 136\% aproximadamente mas agua que en el 2024.
	\newpage
	\begin{center}
		{\huge \textbf{Justificación Del Problema}}
	\end{center}
	\vspace{1cm} 
	{\large %
	Basándonos en lo anteriormente dicho, si dejamos que este 21\% se pierda, a pesar de que sea solo un 21 de 100 sigue siendo alrededor de 300000 toneladas por cosecha/temporada.
	Tal vez este proyecto pueda ayudar a evitar futuras perdidas e influir algo sobre la situación social de la economía en la agricultura.
	Y este proyecto es versátil ya que no solo sería para la plantación de cebolla, la predicción puede ser para varias regiones y solo se pueden modificar los datos y su debido historial de cambio climático a lo largo de los últimos años.
	Otro problema que vimos ante esta situación seria la falta de recursos o estudios ante la gente que va dedicado, sin caer en estereotipos podríamos decir que no tienen mas que la primaria (como se puede plantear mas adelante) así que encontramos 2 soluciones y otras mas que veremos implantar ante el paso del semestre.
	\newpage
	\begin{center}
		{\huge \textbf{Programación}}
	\end{center}
	\vspace{1cm} 
	{\large %
		Hemos encontrado donde podemos meter programación, para generar un código y un algoritmo para usar diferentes fuentes y así poder crear un modelo de lo mas certero en predicciones de lluvia, usando diferentes fuentes y datos de información, para hacer tablas de estadística y usar los datos de manera eficiente.
		
		En esta parte de la programación logramos encontrar una posible solución ante el problema, o bueno sería a quien va dedicado el problema.
		
		Decidimos por parte nuestra iniciar 2 proyectos, pueden servir para ambas secciones sociales ya sea cultivos pequeños como cultivos en masa, como en cultivos mas pequeños.
		
		Uno de estos proyectos justamente es la creación de una app que se actualiza conforme los días y pueda predecir cambios climáticos graves , y una cualidad de alerta por posibles desastres , ya sea huracán , fuertes vientos , inundaciones , lluvias severas.
		
		El otro proyecto es prácticamente lo mismo pero hecho en una pagina , así se puede reducir el espacio de almacenamiento y puede irse actualizando con un simple botón (f5)\
		
		También estamos viendo un posible sistema de riego de bajo consumo y métodos ahorrativos de agua en tiempo de sequía y almacenamiento de agua ante  sobre exceso de lluvias , para su futuro uso 
		\newpage
		\begin{center}
			{\huge \textbf{Ciencia de Datos}}
		\end{center}
		\vspace{1cm} 
		{\large %
		Para esta uca lo que hicimos fue consultar varias fuentes un poco mas directas aunque nos falta tener información de primera mano , para ello haremos unas encuestas en linea y preguntas sobre ámbito social para poder comprender toda el área de investigación , pasarlas por los debidos filtros y así obtener mejores resultados , algunas de las fuentes consultadas estarán cerca del final del documento
		
		En esta uca aprendimos a investigar mas allá de lo que la simple vista puede ser , dejar de ver y empezar a observar , el cuestiona-miento e todo y eso nos ha ayudado a pulir nuestra forma de investigación , esto ira mejorando con el tiempo.
		
		Justamente esta área nos ayudaría a varios temas ,ya sea como en el área de perspectiva de genero , ya que casi se debe de ver toda el área , varios aspectos delicados y mas profundos de los que tienen que verse.
			\newpage
		\begin{center}
			{\huge \textbf{Matematcas Discretas}}
		\end{center}
		\vspace{1cm} 
		{\large %
			Para esta materia nos ha costado encontrar donde meterla, pero podríamos usar las tablas de la verdad junto con la programación para lograr  unos datos mas certeros. 
			
			Aun nos falta encontrar bien donde podríamos meter este lenguaje.
			
			Con las nuevas habilidades que estamos aprendiendo con lo del universo y similares esto podría ayudarnos a encontrar una secuencia o promedio debido al cambio climático si llegamos a lograr el código correcto.
			
		\newpage
	\begin{center}
		{\huge \textbf{Administración Empresarial}}
	\end{center}
	\vspace{1cm} 
	{\large %
			Esta materia la podríamos meter justamente en la organización y el método de creación de este proyecto justamente usando lo que se nos ha enseñado en clase, ya sea dirección, toma de decisiones, el ámbito humano y más factores para la decisión y rumbo de este proyecto,
			
			En el ámbito de sembradío no podríamos influir mucho ya que nos falta experiencia y no somos quienes dirigen el sembradío, pero desde nuestra parte con la creación de proyecto se puede recomendar acciones y unas que otras decisiones.
			“PONEMOS LAS CARTAS EN LA MESA, ELLOS ESCOGEN CUAL CARTA JUGAR”
			Investigando un poco mas sobre el ámbito de la administración empresarial encontramos una información algo preocupante debido a la crudeza de los datos, esto involucra la materia de perspectiva de género ya que en población se usa la mano de obra barata llegando también a la explotación infantil , ya sea por necesidad o  tradiciones familiares .
			
			Cambien vimos que en la área de la administración abarca varias áreas , tanto humanística  como en administración general m enfocándonos un poco mas sobre el área de administración en la siembra de cebolla ya que como es una plantación que surge de la tierra , literalmente se extrae de la tierra el proceso de administración debe ser un poco mas riguroso , ya sea con días y varios procesos de administración .
			
			Justamente un punto de la administración es el valioso recurso humano , ya que la perspectiva de genero se encargara de ver las áreas mas espinosas del tema aquí veremos un tema mas general , siendo uno de los principales el de RECURSOS HUMANOS , este ámbito abarcara el tema muy humanos, ya sea con sus cuidados , entrega de vestimenta reglamentada de protección ,ya sea guantes, casco ,tenis con casquillo, y mas material que puede ser de mucha ayuda ante el trabajador/ra,
			
			La administración empresarial también abarca varios temas , uno importante es la optimizan del trabajo , sea con nuevos métodos de administración, motivación a l@s trabajadorxs, nuevas formas de administración de recursos y mejoras en venta, calidad, trabajo y hasta cooperación y motivación entre los mismos trabajadores para lograr mejores resultados , incluso con compensaciones salariales y de horarios para tener mejor rendimiento y salud de los trabajadores . 
			
		\newpage
	\begin{center}
		{\huge \textbf{Perspectiva de Genero}}
	\end{center}
	\vspace{1cm} 
	{\large %
			En este ámbito se puede usar de una manera más humanitaria, un compañero nos brindó información sobre avance de las mujeres en ámbito de trabajos de cosecha, junto con estadísticas de avance y varias novedades más.
			La información que conseguimos es la siguiente:
			En información general encontramos que el 81\% de la mano e obra es de puros hombres, la cual el 25\% de estos conforman la edad entre 55 y 65 años de edad, el restante ronda entre hombres desde los 18 años hasta los 45
			La mayoría de las regiones grandes sufren de un abuso enorme por parte de la gente que hace la labor pesada (plantación, recolección y más ámbitos) ya que hay testimonios de abuso, explotación infantil y que la mayoría de niñ@s son sacados de la escuela para ir a trabajar por mano de obra barata, esto rondando a los 100 a 150 pesos a cambio de 10-14 hrs de trabajo.
			También hay varios factores , dándose así casos donde los padres con tal de sacar dinero a costa de otros mandan a los niños a trabajar para sacarles el dinero ya que un niño gana prácticamente lo mismo y les quitan el dinero , ya entrando en un caso mas profundo y crudo como la explotación infantil, y esto son solo los casos que llegan a ser entrevistados y ya hasta están normalizados en términos familiares “mi abuelo , mi padre trabajaron desde los 10 años y te toca a ti, salte de la escuela y ponte a trabajar” 
			
			Basándonos en otros datos de varias paginas logramos ver que en promedio de trabajadores hombres seria un total de 88\% con un salario promedio de $3.03k(siendo muy poco debido a la “chinga”) ante la diferencia de un 12\% en mujeres con un salario promedio de $3.49k  y esto por ambas partes es por trabajar alrededor de 37.7hrs a la semana , esta información se encuentra en una pagina oficial en México  aun así nos falta investigar algo ms profundo , coloquial mente se diría hacer “investigación de calle” haciendo encuestas y mas cositas.
			
			\newpage
			\begin{center}
				{\huge \textbf{POSIBLES SOLUCIONES TEMPRANAS}}
			\end{center}
			\vspace{1cm} 
			{\large %
			Hemos encontrado una posible solución temprana junto con varias ideas, las mencionaremos en la siguiente lista:
			1-	Creación de una app, hemos visto mucha técnica, estadística y datos variados, pero hay que recordar que este proyecto va justamente para personas (sin querer sonar a estereotipo) solamente hayan terminado la primaria, por eso la creación de la app de forma sencilla ya que no creemos que la gente tenga una super pc para filtrar semejantes datos.
			
			1-1	chat de voz: vimos la edad general de las personas que cosechan ya sea por varios motivos dejan tanto de leer bien y escuchar bien, así que la app serviría muy bien un chat de vos sencillo para que logren escuchar las estadísticas de manera sencilla y clara 
			2-	Creación de sistema de riego y cultivo: para esto tuvimos la idea de crear un sistema de riego sencillo, usando vaporizadores ya que estos no consumen mucha energía y prácticamente vaporizan el agua , esto lo implementaríamos con un sistema de ventiladores para poder extraer la humedad de manera equitativa y no se sobre humedezcan las zonas donde se genera el vapor
			
			
			3-	Sistema de almacenamiento de agua ante brutales lluvias:
			Aquí aun veremos ideas ya que al ser una cosecha varia la lluvia y agua, ya sea por sequia natural o falta de agua en varias zonas del país , así se pueda filtrar , vitamina y purificar el agua para su uso , el sistema de humidifica-dores lo vimos bien ya que la humedad y puede ayudar a las plantas con vitaminas varias , pesticidas y fertilizantes además de proteger contra el golpe directo de sol y mas cuando el clima va variando tan bruscamente .
			
			4-	La creación de la aplicación, con esto podríamos ayudar ya sea consultando varias fuentes e información y que la aplicación valla variando ante el clima, humedad y estadística ante el paso de los días, combinado del almacén de agua, humidificares, y el sistema de ventilación puede servir en caso de fallar y no se haya podido sembrar (veremos usar comando para poderlo unir a la aplicación y activar el sistema de riego con un solo botón).
			5-	Creación de una pagina , así para abarcar la sociedad que ha sido dejada de lado por la tecnología veremos crear una la p para las producciones que lleguen a tener mas recursos y presupuestos 
	
	\newpage
	\begin{center}
		{\huge \textbf{FUENTES CONSULTADAS}}
	\end{center}
	\vspace{1cm} 
	{\large %
		Las fuentes consultadas son varias, tanto en pdf, como experiencias personales, vídeos descriptivos, estos últimos han sido de mucha ayuda ya que encontramos una serie de vídeos que enseñan desde la siembra de semilla, cortar los primeros brotes y como sembrarlos, métodos de preparación de tierras, como se cultivan, cosechan limpian y venden.
		Hablan también desde perspectivas, ya sea la venta, problemas sociales desde las perspectivas de los cultivadores, precios y acciones que se han tomado, junto con pérdidas y ganancias. 
		-Canal de TIBSITA YOTUBE, muestra desde principio a fin, sembrar semillas, tiempo de cosecha, proceso de limpieza y se comentan situaciones sociales de precios y todo eso 
		https://www.youtube.com/@TISBITA/playlists
		-UNIVERSIDAD DE SONORA: CULTIVO DE CEBOLLA PDF
		-Cesar navarro (compañero que siembra cebolla en casa)
		-Agro climatología y Agrometeorología pdf
		-GUIA PARA PRODUCCION DE CEBOLLA EN ZACATECAS 
		- CONAGUA resúmenes mensuales de temperatura y lluvias
		-gobierno del estado de chihuahua, prensa y reportes 
		-INECC cambio climático en México y la vulnerabilidad de chihuahua ante ese cambio 
		-WEATHER SPARK histórico detallado del clima en chihuahua
		CAACS(UNAM) estado y perspectivas del cambio climático.
	\end{document}

